\section{The problem, results, and proof strategy}\label{ms:introduction}
\subsection{The problem and the candidate constant}
Let $\mathcal P$ be the laws on $\R$ with mean zero, variance one and
finite third absolute moment $\rho(P)=\E|X|^3$. For independent copies
$X_1,\ldots,X_n$ of $P$, put $S_n=\sum_iX_i$ and
\[
 \Delta_n(P)=\sup_{H\in\R}|\Pp_P(S_n\le H)-\Phi(H/\sqrt n)|,
 \qquad C_n=\sup_{P\in\mathcal P}\frac{\sqrt n\,\Delta_n(P)}{\rho(P)}.
\]
Here $\Phi$ and $\phi$ are the standard normal distribution function
and density. The sharp-constant conjecture is
\begin{equation}\label{ms:conjecture}
 \sup_{n\ge1}C_n=\CE,
 \qquad \CE=\frac{3+\sqrt{10}}{6\sqrt{2\pi}}=0.409732\ldots.
\end{equation}
Throughout, $H$ is a threshold for the unnormalized sum $S_n$, and
a terminating decimal in an inequality denotes the corresponding exact
rational number. Affine standardization gives the equivalent result
for any mean and positive variance, with third absolute central moment
divided by the cube of the standard deviation.

Berry and Esseen established the classical order $n^{-1/2}$
\cite{Berry1941,Esseen1942,Esseen1945}. Esseen's lower bound
\cite{Esseen1956} explains the proposed constant. If
$p_*=(4-\sqrt{10})/2$, $q_*=1-p_*$ and
$X=(B-p_*)/\sqrt{p_*q_*}$ for $B\sim\operatorname{Bernoulli}(p_*)$,
the fixed-law lattice limit gives
\begin{equation}\label{ms:bernoulli-limit}
 \lim_{n\to\infty}\frac{\sqrt n\Delta_n(P)}{\rho(P)}
 =\frac{3+q_*-p_*}{6\sqrt{2\pi}(p_*^2+q_*^2)}=\CE.
\end{equation}
Indeed $p_*^2+q_*^2=10-3\sqrt{10}$ and
$(3+q_*-p_*)/(p_*^2+q_*^2)=3+\sqrt{10}$.
Every constant below $\CE$ therefore fails for a sufficiently large sum
of this one law. The upper conjecture is associated with Zolotarev
\cite{Zolotarev1966,Zolotarev1967}; Schulz proves the strict finite-index
bound for every nondegenerate Bernoulli law \cite{Schulz2016}.


\subsection{Results and their hypotheses}\label{ms:scope}
Theorem~\ref{short:global} proves $C_n\le\CE$ for every $n\ge1$,
with strict inequality for $n\ge700$. The Bernoulli limit above proves
that the global constant cannot be lowered. The proof uses three parts.
\begin{center}\small
\begin{tabular}{@{}P{.25\linewidth}P{.68\linewidth}@{}}
\toprule
Part & Conclusion and role \\
\midrule
Finite induction through $42$ &
The first four indices use polynomial majorants and a physical comparison.
For $5\le n\le42$, actual-index support and Fourier bounds acquire the
hypotheses of finite count comparisons. Each step uses sharp earlier
indices and retains one common marginal throughout. \\
First-failure comparison, Theorem~\ref{ext:bridge-composition} &
At a least index with $C_n>\CE$, a one-step smoothing estimate gives
$n(C_n-\CE)<0.3$, sharp deleted predecessors and zero predecessor loss.
A finite comparison excludes $43\le n\le48$; a uniform comparison
excludes every $n\ge49$. This proves the weak bound for all $n$. \\
Compact strict comparison, Theorem~\ref{short:acquired-exclusion} &
For $n\ge700$, a fixed-target maximizer with nonnegative excess would
have $nD\le32.20665378$ and $|H|\le1.83605680$. These bounds place
it in an explicit conditional-support region where the binomial
comparison is strict, including the zero-noise case. \\
\bottomrule
\end{tabular}
\end{center}
The small-index and localization calculations are computer-assisted.
Their proofs combine continuous analytic bounds, exact inequalities on
closed parameter regions, and checks that successive bounds concern
the same law. Appendix~\ref{revision:reproducibility} identifies the
inputs, complete executions and acceptance procedures. The published
structural and finite-binomial results used in the compact comparison
are stated at \eqref{v9:entry:structural} and
\eqref{short:count:binomial}, with their hypotheses and references.

There is no assumption of a sharp bound through $699$. The finite
induction stops at $42$; the first-failure argument supplies all later
indices at once. The separate compact comparison proves strictness
from $700$ onward. Nor is eventual two-point optimality asserted:
the proof bounds every possible extremal witness without replacing
an unrestricted law by a two-point law.

\subsection{Why the defect controls a maximizer}
At a fixed $n$, maximize the difference between the two sides of the
candidate inequality:
\begin{equation}\label{v9:intro:excess}
 a(P,H)=\sqrt n\{\Pp_P(S_n\le H)-\Phi(H/\sqrt n)\}-\CE\rho(P).
\end{equation}
Proposition~\ref{v9:core:attainment} proves attainment. Its first-order
condition bounds $y\mapsto\Pp(S_{n-1}\le H-y)$ by an explicit
piecewise cubic function that agrees with it on the closed support.
This contact condition allows us to use the geometry of that support.

Let $\tau(t)=\E[X\ind_{\{X\ge t\}}]$ be the zero-bias density and
$J=\sum_g\tau_g|g|^2$, where $g$ runs over bounded support gaps and
$\tau_g$ is the constant value of $\tau$ there. A lattice law of span
$h$ has $J\ge h$; if every successive gap has length $h$, then $J=h$.
Thus $J$ measures gap size with zero-bias weights. The inequality
\[
 \gamma+3J\le(3+\sqrt{10})\rho,\qquad \gamma=\E X^3,
\]
proved in Lemma~\ref{ms:physical-defect}, extends the moment--span
bound from lattice laws to arbitrary standardized laws. Its slack
\[
 D=(3+\sqrt{10})\rho-\gamma-3J
\]
is nonnegative and vanishes exactly at the positive-skew Bernoulli
law in \eqref{ms:bernoulli-limit}. The lattice consequence and its
reflection are proved in Corollary~\ref{opus:lattice-gap}.

Integrating the contact relation against $\tau$ in two ways gives the
conservation identity of Theorem~\ref{revision:contact-conservation}.
Since $\CE=(3+\sqrt{10})\phi(0)/6$, its two leading moment terms
combine into a multiple of $D$. Solving for the excess gives
\[
 a=\frac{\mathcal R_c-n\CE\rho(1-e)
             -n\phi(0)eD/6-n^{3/2}L}{n+H^2},
 \qquad e=\exp(-H^2/(2n)).
\]
Here $L\ge0$ is a weighted contribution from the law of $H-S_{n-1}$
inside the support gaps, and $\mathcal R_c$ is an explicitly defined
remainder. The three subtracted terms are nonnegative. A nonnegative
excess therefore requires $n\phi(0)eD/6\le\mathcal R_c$.
Once localization bounds both $\mathcal R_c$ and $|H|$, this bounds
$nD$. The uniform bounds come from Theorem~\ref{v9:entry:theorem},
not from a fixed-law asymptotic expansion.

The predecessor hypothesis supplies an upper bound for
$\Pp(S_{n-1}\le H-t)$ in the contact condition, which excludes remote
support points; see \eqref{v9:entry:contact}. This support argument is
used both in coarse localization and in the final two-ray estimate.
With $nD\le32.20665378$, the geometry at $n\ge700$ gives
$3/8<p<93/200$, inward support displacement below $1/250$ and
$\lambda<7/20$. Count conditioning then compares the law with a
binomial sum. One threshold count controls noise up to $1/100$;
five neighboring counts control the remaining interval. Both
comparisons keep all probabilities entering from other counts.

\begin{figure}[ht]
\centering
\begin{tikzpicture}[node distance=4mm,>=Latex,
 box/.style={draw,rounded corners,align=center,text width=.86\linewidth,
 inner sep=5pt,font=\small}]
\node[box] (contact) {Attainment and contact; nonnegative support-gap defect $D$};
\node[box,below=of contact] (identity) {Conservation and Fourier bounds acquire support, defect and threshold bounds};
\node[box,fill=gray!10,below=of identity] (finite) {Polynomial and count comparisons, with exact closed-domain checks\\
Sharp finite base: $C_n\le\CE$ for $1\le n\le42$};
\node[box,below=of finite] (failure) {A least strict failure has sharp predecessors, zero loss and excess $<0.3/n$};
\node[box,fill=gray!10,below=of failure] (local) {Finite local comparison for $43\le n\le48$; uniform local comparison for $n\ge49$\\
Every possible first failure is excluded};
\node[box,line width=.8pt,below=of local] (global) {$C_n\le\CE$ for all $n$; the compact comparison also gives $C_n<\CE$ for $n\ge700$};
\draw[->] (contact)--(identity);
\draw[->] (identity)--(finite);
\draw[->] (finite)--(failure);
\draw[->] (failure)--(local);
\draw[->] (local)--(global);
\end{tikzpicture}
\caption{The proof's main dependencies. The shaded comparisons combine
analytic inequalities with complete exact finite checks. The uniform
comparison is proved for every index in its range; it is not an
extrapolation from the finite calculation.}
\end{figure}

\FloatBarrier
\subsection{What the certificates prove}
A certificate is finite exact data for a specified mathematical inequality,
together with rules for checking those data. Here the claim is always about
an entire region of possible laws or parameters. It is not a claim about
only the distributions or points tried during a numerical search.

The proof uses three forms of certificate. For the first few indices,
rational polynomial coefficients majorize the closed sum event on a whole
support product; independence expresses the expectation through moments of
one common marginal. For localization and Fourier bounds, rational boxes
cover the complete admissible parameter region. A box is either excluded
by a necessary inequality or assigned bounds valid everywhere on it.
For the count comparisons, a finite tree subdivides the acquired probability,
phase and noise region until each terminal cell has a strict score bound
or a proved contradiction. The checker verifies the initial domains,
every subdivision and all terminal inequalities, including shared
boundaries. Thus every possible witness is accounted for.

Finding these data and verifying them are separate tasks. A search can
choose polynomial coefficients, subdivisions, prices or refinement passes.
It may use floating-point calculations to suggest those choices. Acceptance
uses the resulting exact rational data and checks every claimed inequality;
a successful search or a saved success message has no independent authority.
Some of the present candidates are supplied objects whose complete
acceptance is repeated. Other tables are constructed and then recomputed.
The reproduction record distinguishes these cases. Where the constructor
and acceptor share numerical routines, this is described as exact replay,
not as an independent implementation of the same calculation.

Acceptance uses integer and rational arithmetic. Roots, exponentials,
trigonometric functions and Gaussian probabilities are enclosed by rational
bounds with explicit analytic remainders. The main checking programs use
Python integers; the three-summand Bernstein checker uses C++ with
Boost's arbitrary-precision integers. No floating-point approximation is
accepted as an exact sign test. The trusted implementation includes the
checking programs, their arithmetic libraries and their execution environment;
this is a computer-assisted proof, not a proof-assistant formalization.

A passed certificate still needs its mathematical hypotheses. For example,
a count certificate can establish a bound for every law in a prescribed
support and moment region, while a separate localization argument proves
that the same hypothetical maximizing law lies in that region. The
resource-to-count checks verify this implication without changing the
index, the marginal, or the discrepancy orientation. Continuous arguments
such as attainment, contact, Fourier inversion and the least-failure
induction are proved in the text; they do not follow from an arithmetic
success flag.

The appendices give the inequalities implemented by the checkers and the
rules for their domain coverage. Appendix~\ref{revision:reproducibility}
maps the full argument to its numerical families and explains the source
package, commands and execution records. An independent verification needs these derivations,
the exact inputs, the complete acceptors and the checks matching each
output to the next theorem's hypotheses. Hashes identify the files used;
they do not prove the inequalities.

Subdivision and rigorous enclosure have established precedents in
computer-assisted proofs \cite{SolovyevHales2013}. Numerical optimization
also appears in Berry--Esseen estimates \cite[Section~2]{Shevtsova2011}.
The custom JSON files organize this proof's exact data; their format is
not a standard theorem-proving language. The mathematical justification
comes from the inequalities, complete coverage and correct implementation
of the checking rules.

\subsection{Related work}
Fourier and structural moment inequalities progressively improved the
universal upper bound \cite{vanBeek1972,Prawitz1972,Tyurin2010,
Shevtsova2011,Shevtsova2013}. We use Shevtsova's structural inequality
explicitly in localization. Recent manuscripts report smaller upper
constants \cite{xiao2026berryesseen,lin2026berryesseen}, and He and Cheng
report eventual sharpness beyond a uniform, very large threshold
\cite{hecheng2026eventually}. Work on independent arrays and extremizers
addresses the related non-iid problem \cite{BellIndependent2026,
BellExtremizers2026}. These recent reported results are not used in
our proof. An earlier AI-assisted working draft is on Hexagon
\cite{rodziewiczkolodziejek2026momentdefects}.
